\documentclass[10pt,a4paper]{amsart}
\usepackage[margin=19mm]{geometry}
\usepackage{amsmath,amssymb,mathtools}
\usepackage[scr=rsfs,cal=euler]{mathalfa}
\usepackage{xfrac}
\usepackage[unicode,hidelinks,bookmarks=false]{hyperref}
\newcommand{\ie}{i.e.\ }
\newcommand{\cf}{cf.\ }
\newcommand{\resp}{respectively\ }
\newcommand{\Subsection}{\subsection}
\providecommand{\nobreakdash}{\nobreak\hbox{-}}
\setlength{\emergencystretch}{2em}
\multlinegap0pt
\usepackage{bm}
\usepackage{bbm}
\usepackage{dsfont}
\usepackage{xspace}
\usepackage{cancel}
\usepackage{mathtools}
\usepackage{amsthm}
\makeatletter
\@ifundefined{theorem}{\newtheorem{theorem}{Theorem}}{}
\@ifundefined{lemma}{\newtheorem{lemma}[theorem]{Lemma}}{}
\makeatother
\title[Revisiting local regression: shape regularity, uniform rates]{Revisiting local regression: shape regularity, uniform rates, and the limits of random splits}
\author{J\'er\'emy Bettinger, Fran\c{c}ois Portier, Adrien Saumard}
\date{Source: arXiv:2606.28641v1 (CC BY 4.0)}
\begin{document}
\maketitle

\noindent\emph{Source and licence.} Revisiting local regression: shape regularity, uniform rates, and the limits of random splits. Version arXiv:2606.28641v1 (CC BY 4.0), distributed under the Creative Commons Attribution 4.0 International licence, as recorded in the arXiv licence metadata for this exact version and shown on the abstract page https://arxiv.org/abs/2606.28641v1. Full rights notice: licenses/arxiv-2606.28641-CC-BY-4.0.txt.\par\smallskip

\noindent\textit{Editorial scope.} Counted unit: the Lemma whose source label is \texttt{lemmaWager\_new} in the article (source0.tex lines 1529--1538, source label \texttt{lemmaWager\_new}), delivered together with the article's own complete original proof of it (source0.tex lines 1539--1569, one \texttt{proof} environment of 31 lines). No mathematics is written, repaired or re-derived: statement and proof are the article's own text line for line. Required context retained uncounted: lemmaWager0 (lines 1467--1471). Every definition, display and label of those lines is retained unchanged. Omitted from this package and not redistributed: 1--1466, 1472--1528, 1570--1648. Nothing inside a retained excerpt is omitted or altered: every retained body is delivered exactly as the article's lines are written, so no omission receipt of any kind accompanies this package. Citations: the retained excerpts carry no citation command at all, so no bibliography entry of the article is transcribed and no reference is redistributed. Reference adaptation: none was needed and none was made; every reference inside the delivered unit resolves inside it, so no unresolved reference or citation is produced. Printed slips kept literal: no printed word, symbol, hypothesis or number of the counted statement or of its proof was altered. Layout and class: the article's own class is \texttt{scrartcl} with packages including amsmath, amsfonts, amsthm, bm, dsfont, graphicx, enumitem, natbib, url, colortbl, booktabs, tabularx, float, multirow; the wrapper is the lane's pinned \texttt{amsart} editorial wrapper at 10pt on A4, which restates the theorem-like environments the retained text needs and adds the article's own macro definitions that the retained text needs (none), with extra wrapper packages (bm, bbm, dsfont, xspace, cancel, mathtools). The delivered numbering is the unit's own. Assets: no retained span contains a figure environment, a graphics-inclusion command, a PDF-page inclusion, a table or a TikZ picture; no image file is copied into this package, the package declares no asset dependency and no image file is redistributed with it. Licence: version arXiv:2606.28641v1 of this article is distributed under the Creative Commons Attribution 4.0 International licence, as recorded in the arXiv licence metadata for this exact version and shown on the abstract page https://arxiv.org/abs/2606.28641v1. A full rights notice is delivered at licenses/arxiv-2606.28641-CC-BY-4.0.txt. Characters: every retained excerpt is pure ASCII, so the delivered engine, XeLaTeX with its default Latin Modern text font, reports no missing character.
\begin{lemma}\label{lemmaWager0}
Consider a tree of depth $N$ on $[0,1]^d$. For $k \in  \{0,1,\dots,N \}$, let $\mathcal{V}_k(x)$ be the unique cell of depth $k$ containing $x$. Assume that we have the $( P^X_n,\alpha)$-regularity condition: for all $x\in [0,1]^d$ and  $k \in \{1,\ldots,  N\} $, $$  P_n^X(\mathcal{V}_{k}(x))  \ge \alpha P_n^X(\mathcal{V}_{k-1}(x)).$$ 
Under the condition $16\log( 4 (2n+1)^{2d} / \delta) \leq n \alpha^{N}$, with probability at least $1 - 2\delta$,  we have, for all $x\in [0,1]^d$ and $k \in  \{0,\dots,N \}$, 
$$  \frac{1}{ 4}    P ^X (\mathcal{V}_{k}(x))  \leq  P^X_n(\mathcal{V}_{k}(x))\leq 2 P^X (\mathcal{V}_{k}(x)) . $$
\end{lemma}

\begin{lemma}\label{lemmaWager_new}
Consider a tree of depth $N$ on $[0,1]^d$. For $k \in  \{0,1,\dots,N \}$, let $\mathcal{V}_k(x)$ be the unique cell of depth $k$ containing $x$.
 Assume that we have the $(\alpha,P_n^X)$-regularity condition: for all $x\in [0,1]^d$ and  $k \in \{1,\ldots,  N\} $, 
$  P_n^X(\mathcal{V}_{k}(x))  \ge \alpha P_n^X(\mathcal{V}_{k-1}(x))$. Under the condition $16\log( 4 (2n+1)^{2d} / \delta) \leq n \alpha^{N}$, with probability at least $1 - 2\delta$,  we have, for all $x\in [0,1]^d$ and $k \in  \{1,\dots,N \}$, 
\begin{equation*}
 P^X(\mathcal{V}_{k}(x))  \ge \frac{\alpha}{8} P^X(\mathcal{V}_{k-1}(x)).    
\end{equation*}
Moreover, if $X$ admits a density $f_X$ bounded below by a constant $b > 0$ and above by a constant $M > 0$, then, with probability $1-2\delta$,  we have, for all $x\in [0,1]^d$ and $k \in  \{1,\dots,N \}$, 
$$  \lambda(\mathcal{V}_{k}(x)) \geq \frac{\alpha b}{8M} \lambda(\mathcal{V}_{k-1}(x)) .$$
\end{lemma}

\begin{proof}
For all $k\in \{1,\ldots, N\}$ and all $x\in [0,1]^d$, we have that $P_n^X(\mathcal V_k(x)) \geq \alpha P_n^X(\mathcal{V}_{k-1}(x))$. 
Applying Lemma \ref{lemmaWager0}, we obtain with probability $1-2\delta $, 
for all $k\in \{1,\ldots, N\}$ and all $x\in [0,1]^d$,
$$ 2 P^X(\mathcal V_k(x)) \geq P_n^X(\mathcal V_k(x)) \geq \alpha P_n^X(\mathcal{V}_{k-1}(x)) \geq \frac{\alpha }{4} P^X(\mathcal V_{k-1}(x)).$$
Using the boundedness assumptions on $f_X$, we obtain, for any hyper-rectangle $V\subset [0,1]^d$, $M \lambda(V) \geq  P^X(V ) \geq b \lambda(V)$, which allows to conclude.
%P_n^X(\mathcal V_k(x))  \geq  \alpha^{k} P_n^X(\mathcal{V}_0(x)) = \alpha^{k}\geq \alpha^N $. As a consequence, our condition $16 \log( 4 (2n+1)^{2d} / \delta) \leq n \alpha^{N}$ implies that 
% $$ \text{for all }x\in [0,1]^d \text{ and } k\in \{0,1,\ldots, N\} , \qquad \frac{ 4 \log( 4 (2n+1)^{2d} / \delta) }{ n P_n^X(\mathcal V_k(x)) } \leq \frac 1  4. $$ 
% Using Vapnik's inequality (last statement in Theorem \ref{th:vapnik_normalized}) with the class $\mathcal{A}$ of hyper-rectangles in $\mathbb{R}^d$ because $\mathcal{V}^k \subset \mathcal{A}$, we have $\mathcal{S}_{2n}(\mathcal{A}) \leq (2n+1)^{2d}$. As a consequence, with probability $1-\delta$, for all hyper-rectangles $V$,
% \begin{align*}
%     P^X(V) &\geq P_n^X (V) \left(1-\sqrt { \frac{4 \log( 4 \mathcal{S}_{2n}(\mathcal{A}) / \delta) }{ nP_n^X (V)} } \right) \\
%     &\geq P_n^X (V) \left(1-\sqrt { \frac{4 \log( 4 (2n+1)^{2d}  / \delta) }{ nP_n^X (V)} } \right)\geq \frac 1 2 P_n^X (V) 
% \end{align*}
% The latter being valid for all hyper-rectangles, it must be true for $V = \mathcal{V}_{k}(x) $, for all $k = 0,1,\ldots, N$ and all $x\in [0,1]^d$, leading to, with probability $1-\delta$, for all  $x\in [0,1]^d$ and all $k\in \{ 0,\ldots, N\}$,
% \begin{align}\label{event1}
%   P^X (\mathcal{V}_{k}(x) ) \geq \frac 1 2  P_n^X (\mathcal{V}_{k}(x) ),%\qquad   P^X (\mathcal{V}_k(x) \backslash \mathcal{V}_{k+1}(x)\} ) \geq \frac 1 2  P_n^X ( \mathcal{V}_k(x) \backslash \mathcal{V}_{k+1}(x) )    
% \end{align}
% % $$\min (P^X (\mathcal{V}_{k+1}(x) )  , P^X (\mathcal{V}_k(x) \backslash \mathcal{V}_{k+1}(x)\}  )  \geq \frac 1 2 \alpha P_n^X (\mathcal{V}_k(x))    .$$
% Note that the above inequality implies that for all $x\in [0,1]^d$ and $k\in \{0,1,\ldots, N \}$,
% $$ \frac{ 4 \log( 4 (2n+1)^{2d} / \delta) }{ n P^X(\mathcal V_k(x)) } \leq \frac{ 8 \log( 4 (2n+1)^{2d} / \delta) }{ n P^X_n(\mathcal V_k(x)) }\leq \frac 1  2. $$ 
% In a similar way as before, we now apply another Vapnik's inequality (first statement in Theorem \ref{th:vapnik_normalized}) to obtain that, with probability at least $1 - \delta$, for  all $x\in [0,1]^d$ and all integer $k \in \{0,\dots,N\}$,
% \begin{align}\label{event2} 
% P^X_n(\mathcal{V}_{k}(x))  &\geq P ^X (\mathcal{V}_{k}(x)) \left(1-\sqrt { \frac{4 \log( 4 (2n+1)^{2d} / \delta) }{ nP ^X (\mathcal{V}_{k}(x))  } } \right) 
% \end{align}
% Combining both events \eqref{event1}, \eqref{event2} and the $\alpha$-regularity condition, it follows that, with probability $1-2\delta$, for all $x\in [0,1]^d$,and  $k\in \{ 0,\ldots , N-1\}$,
% $$ P^X (\mathcal{V}_{k+1}(x)  \geq \frac 1 2 \alpha P^X_n(\mathcal{V}_{k}(x)) \geq \frac 1 2 \left(1-\sqrt { \frac{1}{ 2} } \right)  \alpha P ^X (\mathcal{V}_{k}(x))   . $$
% % and
% % $$ P^X ( \mathcal{V}_k(x) \backslash \mathcal{V}_{k+1}(x)\})  \geq \frac 1 2 \alpha P^X_n(\mathcal{V}_{k}(x)) \geq \frac 1 2 \left(1-\sqrt { \frac{1}{ 2} } \right)  \alpha P ^X (\mathcal{V}_{k}(x))   ,  $$
% Noticing that $ (1- \sqrt {1 / 2} ) \geq 1/4$ allows to conclude.

\end{proof}

\end{document}
